% Хэрэглээний тохиолдлын (use case) диаграмм (TikZ).
% Эх: demo_server.py-ийн нээлттэй API (/api/route, /api/voice/transcribe, /api/voice/speak,
% /api/admin/*) ба tools/import_school_tables.py. Оролцогчдыг 2026.11.12-ны ресепшний
% мэдээллээр хоёр бүлэгт хуваасан: өдөр тутмын зочин, нээлттэй өдөрлөгийн зочин.
% Байрлалыг (x,y) координатаар өгсөн.
\begin{tikzpicture}[
    font=\footnotesize,
    uc/.style={draw, ellipse, align=center, text width=3.3cm, minimum height=0.95cm, inner sep=1pt},
    actor/.style={align=center, font=\footnotesize},
    extbox/.style={draw, align=center, text width=2.5cm, minimum height=0.9cm, fill=black!8},
    line/.style={semithick},
    dep/.style={dashed, -{Stealth[length=2mm]}},
    gen/.style={-{Triangle[open, length=2.6mm, width=2.6mm]}, semithick},
    lab/.style={font=\scriptsize, inner sep=1.5pt, fill=white}]

% --- хүний дүрс зурах макро: #1 нэр, #2 (x,y), #3 бичиг
\newcommand{\stickman}[3]{%
  \begin{scope}[shift={(#2)}]
    \draw[semithick] (0,0.55) circle (0.17);
    \draw[semithick] (0,0.38) -- (0,-0.1) (-0.3,0.2) -- (0.3,0.2) (0,-0.1) -- (-0.25,-0.5) (0,-0.1) -- (0.25,-0.5);
    \node[actor, anchor=north] at (0,-0.55) {#3};
    \coordinate (#1) at (0.35,0.1);
    \coordinate (#1-l) at (-0.35,0.1);
  \end{scope}}

% --- системийн хүрээ
\draw[semithick] (2.4,2.1) rectangle (11.0,-9.6);
\node[font=\footnotesize\bfseries] at (6.7,1.75) {Сургуулийн ухаалаг туслах};

% --- хэрэглээний тохиолдлууд
\node[uc] (typed)  at (4.6,0.7)   {Бичгээр асуух};
\node[uc] (voice)  at (4.6,-0.8)  {Дуугаар асуух};
\node[uc] (answer) at (4.6,-2.5)  {Эх сурвалж, огноотой\\хариулт авах};
\node[uc] (refer)  at (4.6,-4.2)  {Хэнд, хаана хандахыг\\мэдэх};
\node[uc] (listen) at (4.6,-5.9)  {Хариултыг дуугаар\\сонсох};
\node[uc] (stt)    at (8.9,-0.8)  {Яриаг текст болгох};
\node[uc] (tts)    at (8.9,-5.9)  {Текстийг дуу болгох};
\node[uc] (import) at (8.9,-7.4)  {Сургуулийн хүснэгт\\оруулах};
\node[uc] (publish) at (8.9,-8.8) {Мэдээллийн шинэ хувилбар\\шалгаж нийтлэх};

% --- оролцогчид
\stickman{guest}{0.4,-2.4}{Зочин}
\stickman{daily}{-3.2,-0.4}{Өдөр тутмын зочин\\{\scriptsize оюутан, эцэг эх,}\\{\scriptsize ажил горилогч}}
\stickman{open}{-3.2,-4.4}{Нээлттэй өдөрлөгийн\\зочин\\{\scriptsize элсэгч, эцэг эх}}
\stickman{staff}{13.4,-8.0}{Сургуулийн\\ажилтан}
\node[extbox] (gemini) at (13.4,-0.8) {Яриа таних\\үйлчилгээ\\{\scriptsize Gemini 3.5 Transcribe}};
\node[extbox] (edge)   at (13.4,-5.9) {Дуу үүсгэх\\үйлчилгээ\\{\scriptsize Edge TTS}};

% --- холбоос
\foreach \u in {typed, voice, answer, refer, listen} \draw[line] (guest) -- (\u.west);
\draw[gen] (daily) -- (guest-l);
\draw[gen] (open) -- (guest-l);
\draw[dep] (voice) -- node[lab] {«include»} (stt);
\draw[dep] (listen) -- node[lab] {«include»} (tts);
\draw[dep] (refer) -- node[lab, right] {«extend»} (answer);
\draw[line] (stt) -- (gemini);
\draw[line] (tts) -- (edge);
\draw[line] (staff-l) -- (import.east);
\draw[line] (staff-l) -- (publish.east);
\end{tikzpicture}
